\chapterphoto{UCSC Baskin 工程楼旁的树木}{photo-chap0}
\chapter{致读者}

\section{预备知识}
本书假定读者具备基本的代数与复数知识。必要的入门物理与微积分知识会在需要时予以讲授。

\section{本书的结构}
\label{sec:structure}
本书由五个部分和一组附录组成，涵盖控制工程师所要完成的四项任务：推导系统的模型（运动学）、为模型设计控制器（控制理论）、设计观测器来估计模型的当前状态（定位），以及规划控制器如何驱动模型到达期望状态（运动规划）。

第一部分``控制理论基础''介绍控制理论的基础知识，讲授 PID 控制器设计的基本原理。

第二部分``现代控制理论''首先以速成课的形式介绍线性代数背后的几何直觉，并讲授足够多的矩阵代数运算方法，使读者能够跟上后续章节的内容。它介绍状态空间表示、能控性与能观性。第一部分获得的直觉与线性代数的记号将被用于对线性多输入多输出（MIMO）系统进行建模和控制，内容涵盖离散化、LQR 控制器设计、LQE 观测器设计和前馈。随后，这些概念被应用于实际系统的控制器设计与实现。第四部分的例子会被转换为状态空间表示，加以实现，并用离散控制器进行测试。

第二部分还介绍了利用李雅普诺夫函数进行非线性控制系统分析的基础知识。它给出了一个独轮车（unicycle）型车辆非线性控制器的例子，以及如何将其应用于两轮车辆。由于非线性控制并不是本书的重点，我们会列出其他书籍和资源供进一步阅读。

第三部分``估计与定位''介绍随机控制理论领域。本书讲授 Luenberger 观测器以及卡尔曼滤波背后的概率论，并给出若干卡尔曼滤波理论创造性应用的例子。

第四部分``系统建模''介绍推导前述章节模型所需的基本微积分和物理概念，并逐步推导若干常见 FRC 子系统的模型。然后讨论系统辨识方法，用于通过实验测量模型参数。

第五部分``运动规划''讨论如何规划机器人从当前状态到达某个期望状态的路径，使运动在其动力学可达的范围内。它介绍用于简单机动的一自由度运动速度曲线（motion profile）。对于自由度更高的曲线，则介绍轨迹优化方法来生成。

附录提供了进一步提高的内容，并非理解本书主体所必需。其中包括对正文中所给出并使用的许多结论的推导。

案例研究中用于生成图像的 Python 脚本，同时也是相应章节所讨论技术的参考实现。它们可以在本书的 Git 仓库中找到，仓库位置见版权页。

\section{本书的理念}
\label{sec:ethos}
本书既可以作为新学生的教程，也可以作为需要温习某些内容的更有经验的读者的参考手册。虽然本书并不追求面面俱到，但希望读者读完之后能够学有所成：要么能亲自实现书中介绍的概念，要么知道去哪里查找更多信息。

书中的某些部分在数学上是严谨的，但我相信，应当让学生打下扎实的理论基础，同时侧重直觉，这样他们才能把所学应用到新的问题上。要深入理解本书的主题，数学是不可避免的。话虽如此，我仍会尽可能提供实用、直观的解释。

大多数教学资源都把线性控制和非线性控制分开讲授，后者通常留给另一门课程。在这里，我们把它们放在一起介绍，因为非线性控制的概念经常出现，而且（如果不考虑李雅普诺夫稳定性的话）这个跨越并不算大。控制与估计各章会在适当的时候介绍处理非线性的相关工具，例如线性化。

\section{无私工程师的心态}
\label{sec:mindset}
工程不只是一套技能，更是一种心态。工程师有一种独特的问题处理方式。下面的格言概括了我希望传授给我的机器人学生们的东西（例子取自控制工程领域）。

\begin{quote}
\itshape ``基于需求做工程设计，而不是基于某种意识形态。''
\end{quote}

工程充满了权衡。工具应当适配任务，并非每个问题都是等着被锤子敲的钉子。正确的做法是：评估解决当前任务所需的最低要求（最低规格），然后只做刚好能满足它们的工作；超出规格就是浪费时间和金钱。如果你需要高于最低规格的性能或可维护性，那么按定义，你的最低规格就选错了。

控制工程在类似的意义上也是务实的：解决。眼前。的问题。对于非线性系统的控制，被控对象求逆（plant inversion）在纸面上很优雅，但在模型不准确时并不奏效；而使用理论上``不正确''的解——比如对非线性系统做线性近似——效果却足够好，以至于整个工业界都在使用它。有比 PID 更复杂的控制器，但我们仍然使用 PID，就是因为它通用而简单。有时，``较差''的解反而比控制系统设计者心目中理想或``干净''的方案更有效，或者具有更划算的收益成本比。请选择最有效的工具。

解决方案只需要足够好，而不必完美。我们想避开积分器，因为它们会引入不稳定，但我们仍然使用它们，因为它们能很好地满足跟踪指标。不应盲目地捍卫某个设计或遵循某种教条，因为总存在其对立面才是更优选择的情形。工程师应当能够判断何时属于这种情况，放下自我，去做能满足客户要求的事情（例如系统响应特性、可维护性、易用性）。出于务实或技术上的原因而偏好某个方案是没有问题的，但工程师不应在个人情感上纠结于最终选择了哪个``足够好''的方案。

\section{征求反馈}
\label{sec:feedback}
虽然我们努力写出一本让控制理论变得平易近人的书，但对某些读者而言，它可能仍然显得内容密集或节奏偏快（这本薄书涵盖了三门反馈控制课程的内容，其中两门是研究生课程）。请通过版权页上列出的 GitHub 链接向我们发送反馈、勘误或建议。也欢迎提供能够展示关键概念、并使它们更易于理解的新例子。
